% Einheit 1 von 2 – Einheit 1 (bank/vierfeldertafel/e1.jsonl, zusammenbau v0.1)
%% TODO Prüfungswort Einheit 1: nicht in der Marken-Zeile
%% TODO Zeitmarke Einheit 1: Marken-Zeile nicht lesbar
%% TODO Zweigzeile Teil 1: was hier gelernt wird (Satz in Schülersprache aus dem Einheitstitel) – Einheit 1
%% TODO Einheitentitel 1 nicht in der Mappe lesbar
\einheitenkopf[e1]{Einheit 1 von 2 · Einheit 1}
\setcounter{aufgabe}{7}

%% TODO Titel: Ich-kann-Satz fehlt in der Bank (Platzhalter: Kettenname) – Nr. 8
%% TODO Anweisung: ein Satz über der ersten Teilaufgabe fehlt in der Bank – Nr. 8
\begin{aufgabe}{Weder–noch, entweder–oder, oder?}
\begin{teile}
\teil „Die befragte Person hat weder einen Hund noch eine Katze.“ Wie viele Felder der Vierfeldertafel sind gemeint? \\ \kreuz{ein Feld} \\ \kreuz{zwei Felder} \\ \kreuz{drei Felder}
\end{teile}
\end{aufgabe}

%% TODO Titel: Ich-kann-Satz fehlt in der Bank (Platzhalter: Kettenname) – Nr. 9
%% TODO Anweisung: ein Satz über der ersten Teilaufgabe fehlt in der Bank – Nr. 9
\begin{aufgabe}{Die Tafel füllen}
%% TODO \rechenplatz{n}: Schrittzahl des Grundfalls fehlt in der Bank (nur bei fester Schreibform, 2.3 b)
\begin{teile}
\teil In einem Sportverein gilt: „$40\,\%$ der Mitglieder spielen Tennis.“ Was ist diese Angabe? \\ \kreuz{Rand} \\ \kreuz{Feld} \\ \kreuz{bedingter Anteil}
\teil In einer Stadt haben $40\,\%$ der Haushalte ein Zeitungsabo ($Z$), $30\,\%$ sind Seniorenhaushalte ($S$), und $18\,\%$ sind Seniorenhaushalte mit Zeitungsabo. Vervollständige die Vierfeldertafel (in Prozent).

\vierfeldertafel{S}{Z}{18,,40,,,,30,,100}
\teil In einer Stadt nutzen $45\,\%$ der Einwohner das Fahrrad als wichtigstes Verkehrsmittel ($R$). $40\,\%$ der Einwohner sind Studierende ($S$); von den Studierenden nutzen $70\,\%$ vor allem das Rad. Stelle den Sachverhalt in einer Vierfeldertafel dar, und mit welcher Wahrscheinlichkeit nutzt eine zufällig ausgewählte Person, die nicht studiert, vor allem das Rad? \hfill (Abitur 2018 LK) \\

\vierfeldertafel{R}{S}{}
\leerfeld[\%]
\teil Bei einem Leihfahrrad ist mit $8\,\%$ Wahrscheinlichkeit die Bremse defekt ($B$), mit $5\,\%$ das Licht ($L$), und mit $89\,\%$ ist weder die Bremse noch das Licht defekt. Vervollständige die Vierfeldertafel (in Prozent). \hfill (Abitur 2018 GK)

\vierfeldertafel{L}{B}{,,8,,89,,5,,100}
\teil Bei einem Tag der offenen Tür stellen $40$ Schülerinnen und $60$ Schüler je eine Frage. $70\,\%$ aller Fragen betreffen das Studium, die übrigen die Wohnsituation. Von den Fragen der Schülerinnen betrifft ein Viertel die Wohnsituation. Wie viele Fragen der Schüler betreffen die Wohnsituation? \hfill (Abitur 2017 LK) \\ \leerfeld
\teil Ein Anbieter hat die Tarife Basis ($30\,\%$ der Kunden) und Premium ($20\,\%$); alle übrigen Kunden haben den Tarif Plus. $55\,\%$ aller Kunden nutzen die App, darunter die Hälfte der Plus-Kunden; $10\,\%$ aller Kunden haben den Tarif Basis und nutzen die App nicht. Vervollständige die Tabelle (in Prozent). \hfill (Abitur 2022 LK)

\sachtabelle{lcccc}{ & Basis & Plus & Premium & Summe}{App & \leerzelle & \leerzelle & \leerzelle & 55 \\ keine App & 10 & \leerzelle & \leerzelle & \leerzelle \\ Summe & 30 & \leerzelle & 20 & 100}
\teil Für zwei Ereignisse gilt $P(A \cap B) = p$, $P(A) = 2p$ und $P(B) = 6p$ mit $p > 0$. Vervollständige die Vierfeldertafel mit Termen in $p$ und zeige, dass $p$ nicht $\frac{1}{6}$ sein kann. \hfill (Abitur 2021 LK)

\vierfeldertafel{B}{A}{p,,2p,,,,6p,,1}
\end{teile}
\end{aufgabe}

%% TODO Titel: Ich-kann-Satz fehlt in der Bank (Platzhalter: Kettenname) – Nr. 10
%% TODO Anweisung: ein Satz über der ersten Teilaufgabe fehlt in der Bank – Nr. 10
\begin{aufgabe}{Die Tafel füllen – Fehler finden, Begründen}
\begin{teile}
\teil Ich finde den Fehler: $50\,\%$ der Kunden eines Ladens sind Stammkunden ($S$); von den Stammkunden zahlen $60\,\%$ mit Karte ($K$). Nele trägt in das Feld $S \cap K$ der Vierfeldertafel $60\,\%$ ein.
\teil Begründe, warum ein Anteil „innerhalb einer Gruppe“ ein bedingter Anteil ist und erst die Multiplikation mit dem Rand das Feld liefert.
\end{teile}
\end{aufgabe}
